\section{Conclusion}
\label{sec:conclusion}

A low-cost AD5941/RP2040 potentiostat was found to recover a resistance
$424\,\%$ in error from cyclic voltammetry---its reported currents low by a
factor of $5.22$---while passing a resistive-dummy-cell linearity check at
$R^{2} > 0.99999$. The cause was a single firmware constant declaring the
value of an on-board calibration resistor that did not match the resistor
fitted, compounded by a broken path from the user-facing command to the code
that consumes it and by a calibration that ran once per power-up. Because a
constant gain error enters the reported current multiplicatively, and because
the coefficient of determination is invariant under affine rescaling of the
measured variable, no linearity statistic could have detected it---and we
verified experimentally that $R^{2}$ did not move measurably while the induced
error was swept over more than two decades.

We proposed a five-check validation protocol that places scale traceability
first, compares recovered values against component nominals rather than against
a fitted slope, reports fit residuals against the converter's quantisation step,
and requires invariance across acquisition settings. It needs no reference
instrument. Applied after correcting eight firmware defects, it placed the
instrument's cyclic voltammetry at $(0.43 \pm 0.02)\,\%$ of the nominal
\SI{10560}{\ohm} reference with a replicate relative standard deviation of
$0.015\,\%$, invariant across scan rate, potential window and potential step.

We also report what the same protocol shows about the instrument's other
methods: chronoamperometry, square-wave and differential pulse voltammetry
return no cell current on this hardware, and impedance spectroscopy does not
complete a frequency point, although the correction of an unbounded wait now
makes that failure fast and diagnosable rather than indefinite. One method is
validated; five are not. That asymmetry, on a single board driven by a single
front end, is the clearest argument we can offer that instrument-level
validation claims should be replaced by per-method ones.

Two recommendations generalise beyond this instrument. First, publish the
calibration: an instrument that reports the gain constants actually in force
turns an entire class of silent, scale-corrupting defects into visible ones, at
the cost of one line of output. Second, report accuracy against an independently
known value, and state the basis of that value; if the only evidence available
is a linearity statistic, say so, and say what it does and does not bound.

Future work should characterise the fitted calibration resistor against a
traceable standard to separate the residual $0.3\,\%$ instrument term, extend
validation across transimpedance ranges and current decades, repair the
low-power loop drive in the direct-current method classes and the impedance
acquisition path, and factor the triplicated configuration and conversion code
into a single shared component so that a defect found once is fixed once.
